\section*{Chapter \ref{ch:iv1:conditional-probability-and-bayes} --- Conditional Probability and Bayes}
\markboth{\protect\omTocMark{}{Answers: Conditional Probability and Bayes}}{\protect\omTocMark{}{Answers: Conditional Probability and Bayes}}

\begin{solution}{iq:iv1:conditional-probability-and-bayes:1}
Told ``at least one is a girl'', the probability is $1/3$. Having met a random child who is a girl, it is $1/2$.

The four families BB, BG, GB, GG are equally likely. ``At least one girl'' removes only BB, and of the three left, one is GG. When you meet a random child, you learn more. Family GG shows you a girl for sure, while BG and GB show you one half of the time, so the weights are $1 : \tfrac12 : \tfrac12$, and GG has probability $1/2$.

\iqwrong{Assuming the two statements are the same fact. They differ in how the information reached you\index{protocol of information}: what matters is the probability of hearing this message under each possibility, not merely whether the message is true.}
\end{solution}

\enlargethispage{-3\baselineskip}%
\begin{solution}{iq:iv1:conditional-probability-and-bayes:2}
The probability is $1/6$, about 17\%.

Use Bayes by odds\index{Bayes by odds}. The prior odds of disease are $1 : 99$. A positive test is $0.99/0.05 = 19.8$ times as likely with the disease as without it, so that is the likelihood ratio\index{likelihood ratio}. The posterior odds are $19.8 : 99 = 1 : 5$, which is a probability of $1/6$.

In people: of 10,000, 100 are ill and 99 of them test positive. Of the 9,900 well, 495 test positive. So 99 of the 594 positives are ill.

\iqwrong{Answering 99\%, or a figure near it. That confuses $P(\text{positive} \mid \text{ill})$ with $P(\text{ill} \mid \text{positive})$, and it ignores how rare the disease is.}
\iqwhy{The same sum sets the value of any alert. A signal that is right 99\% of the time on rare events still produces mostly false alarms.}
\end{solution}

\begin{solution}{iq:iv1:conditional-probability-and-bayes:3}
Switch: you win with probability $2/3$. With 100 doors, you win with probability $99/100$.

Your first pick is right with probability $1/3$. The host's action tells you nothing about your own door, because the host can always open a door with nothing behind it, whichever door you hold. So that door keeps its $1/3$, and the remaining door collects the other $2/3$. Equivalently, switching wins exactly when your first pick was wrong.

The hundred-door version makes this plain: your door is right with probability $1/100$, and the host, who knows, has shown you where the prize would be in all the other cases.

\iqwrong{``Two doors are left, so each has one half.'' Two outcomes do not mean equal chances. The host's knowledge is what makes the update unequal, and a host who opened at random would give a different answer.}
\end{solution}

\begin{solution}{iq:iv1:conditional-probability-and-bayes:4}
The probability is $1024/1123 \approx 0.912$.

By odds\index{Bayes by odds}, the prior odds of the two-headed coin are $1 : 99$. Each head is twice as likely under that coin, so each head doubles the odds. After 10 heads, the odds are $1024 : 99$, which gives $1024/1123$.

The belief crosses one half between the sixth and the seventh head, since $2^6 = 64 < 99 < 128 = 2^7$.

\iqwhy{Odds make this a one-line update and show how quickly repeated evidence overwhelms a prior. Ten doublings multiply by about a thousand, so a prior of 1 in 100 falls fast, but a prior of 1 in a million would not.}
\end{solution}

\begin{solution}{iq:iv1:conditional-probability-and-bayes:5}
The probability is $2/3$.

Count the red sides, not the cards. There are three red sides, and you are equally likely to be looking at any of them. Two of those three belong to the red-red card, and their other side is red. One belongs to the mixed card, and its other side is blue.

\iqwrong{One half, because only two cards could be on the table. That treats the two possible cards as equally likely, but the red-red card is twice as likely to show you a red side.}

\end{solution}

\begin{solution}{iq:iv1:conditional-probability-and-bayes:6}
The probability is $12/29 \approx 0.414$, so the taxi is still more likely to be green.

Use odds\index{Bayes by odds}. The prior odds of blue are $15 : 85$. The witness is 4 times as likely to say blue when it is blue as when it is green ($0.8/0.2$). So the posterior odds are $60 : 85 = 12 : 17$, and that gives $12/29$.

\iqwrong{80\%, the witness's reliability. That answers ``how often is the witness right about any taxi'', not ``given this statement about a rare colour, how often is it right''. Green taxis are so common that their mistaken reports outnumber the true reports about blue ones, 17 to 12.}
\iqfollow{And if the witness were right half of the time?}{Then the likelihood ratio is 1, and the answer stays at the prior, 15\%.}
\end{solution}

\begin{solution}{iq:iv1:conditional-probability-and-bayes:7}
Both six has probability $1/11$, and sum 7 has probability $2/11$.

Of the 36 rolls, 11 contain a six: 6 with the first die a six, 6 with the second die a six, with $(6, 6)$ counted twice. Only one of the 11 is $(6, 6)$, and two of them add up to 7, which are $(1, 6)$ and $(6, 1)$.

\iqwrong{One sixth, by thinking ``one die is a six, so the other is an independent die''. That is right only if you are told the first die is a six. ``At least one'' is not about a particular die, and that changes the count.}
\iqfollow{And if the speaker looks at one die at random and says it is a six?}{The chance of both being six is then $1/6$, as the protocol of information has changed.}
\end{solution}

\begin{solution}{iq:iv1:conditional-probability-and-bayes:8}
The probability is $2/3$.

Say the original ball is red with probability $1/2$. If it is red, the bag holds two reds and the draw is red for certain. If it is blue, the draw is red with chance $1/2$. So by Bayes\index{Bayes by odds} the odds that the original ball was red, given a red draw, are $1 : \tfrac12$, which is $2/3$. If the original ball was red, the ball left in the bag is red; if it was blue, the ball left is blue.

It is also possible to count the balls directly. There are three equally likely reds that could have been the drawn ball: the added red, the original red in the red case, and the added red in the blue case. In two of the three, a red ball remains.

\iqwrong{One half, by arguing that the draw took out the red you added, so the bag is as it was. The red draw is evidence of the original colour, because it is more likely when the bag held two reds.}
\end{solution}

\begin{solution}{iq:iv1:conditional-probability-and-bayes:9}
With a random guard, A stays at $1/3$ and C rises to $2/3$. If the guard always names B whenever possible, A and C are at $1/2$ each.

The likelihood of the guard's words decides everything. If C is pardoned, the guard must say B. If A is pardoned, the guard has a choice, and a random guard says B half the time. So the weights for A and C are $\tfrac13 \cdot \tfrac12$ and $\tfrac13 \cdot 1$, and the odds are $1 : 2$. A's chance has not moved, because the guard could say something like this in any case, and the news moves C's chance.

If the guard always says B whenever possible, then B is certain when A is pardoned, and the weights are equal, so each has $1/2$. And a guard who names C whenever possible gives a stark result: hearing B then means that C is the one pardoned.

\iqwrong{``Two prisoners are left, so one half.'' It holds only under one particular rule for the guard. A probability is fixed not only by the list of remaining possibilities but also by the rule that made the message\index{protocol of information}.}
\iqwhy{When someone tells you a fact selected from what they know, ask how they chose it. This is what separates a filtered report from a neutral sample.}
\end{solution}

\begin{solution}{iq:iv1:conditional-probability-and-bayes:10}
Under the first rule it does not matter: each door has probability $1/2$. Under the second rule, switching wins for certain after door 3 opens, and with chance $1/2$ after door 2 opens.

First rule. The host has no information, and opening an empty door is itself evidence. If the prize is behind door 1, then an empty door is certain; otherwise it has chance $1/2$. So the weights are $\tfrac13 : \tfrac23 \cdot \tfrac12$, which is even.

Second rule. If the prize is behind door 2, the host must open door 3. If it is behind door 1 there is a choice, and the host always opens door 2. So seeing door 3 opened tells you that the prize is behind door 2: switching wins with certainty. Seeing door 2 opened leaves door 1 or door 3, with equal weight, and either of them could have led to it, so the chance is $1/2$.

Averaged over what the host does, door 3 opens a third of the time and door 2 opens two-thirds of the time, so switching still wins $\tfrac13 \cdot 1 + \tfrac23 \cdot \tfrac12 = \tfrac23$, as it must.

\iqwrong{Giving $2/3$ under every rule. The famous answer assumes a knowing host with no preference. A preference does not change the average, but it changes what you can infer from what you actually saw.}
\end{solution}

\begin{solution}{iq:iv1:conditional-probability-and-bayes:11}
The probability is $13/27$.

Count, using 14 types of child, one for each sex and day. There are $14^2 = 196$ equally likely families. Those with at least one Tuesday boy have the Tuesday boy first (14 families) or second (14 families), with the one double-counted family removed: 27. Among those, both children are boys in $7 + 7 - 1 = 13$.

The surprising day matters because a more specific description makes it less likely that two different children could fit it, so it moves you from ``at least one boy'' towards ``this particular boy''. With no day stated, the answer is $1/3$. With a very rare trait, such as a birthday, the answer is $(2 \cdot 365 - 1)/(4 \cdot 365 - 1)$, nearly $1/2$.

\iqwrong{Saying the day is irrelevant and answering $1/3$. The question asks only about the sexes, but a family with two boys has nearly twice the chance ($13/49$ against $7/49$) of containing a Tuesday boy.}
\iqfollow{Is the answer sensitive to how you came to hear this?}{Yes. If a parent picks one child at random and reports that child's sex and day, the answer is $1/2$. Only a rule that says ``report if you have a Tuesday boy'' gives $13/27$.}
\end{solution}

\begin{solution}{iq:iv1:conditional-probability-and-bayes:12}
The probability is $2/3$. It is not $0.7$ because your uniform prior still counts, as if you had seen one head and one tail before you started.

The likelihood of 7 heads in 10 flips with probability $p$ is proportional to $p^7(1 - p)^3$, so the posterior is a Beta distribution\index{Beta posterior} with parameters $8$ and $4$. Its mean, which is the chance of a head next, is $8/12$. In general, Laplace's rule of succession\index{Laplace's rule of succession} gives $(h + 1)/(n + 2)$.

With no data it gives $1/2$, and it never gives 0 or 1 after finitely many flips, which is what you want from a rule used for rare events.

\iqwrong{Using the proportion $0.7$. That is the most likely value of $p$, not its mean, and the posterior is not symmetric around its peak. A heavier prior would shrink further towards one half.}
\iqwhy{Shrinkage towards a prior is built into any estimate like this. A hit rate from few trades should be pulled back towards a base rate before being trusted.}
\end{solution}

\begin{solution}{iq:iv1:conditional-probability-and-bayes:13}
The probability is $35/52 \approx 0.673$.

First update, then predict. The two heads have probability $9/16$ under the biased coin and $1/4$ under the fair one, so by odds\index{Bayes by odds} the biased coin has $9 : 4$, which is $9/13$. The chance of a third head is then the mix\index{law of total probability}
\[
\tfrac{9}{13} \cdot \tfrac34 + \tfrac{4}{13} \cdot \tfrac12 = \tfrac{35}{52}.
\]
Before any flips, this mix would be $5/8$. So each head moves the next forecast up, because you are learning about the coin.

\iqwrong{Giving $5/8$ on the ground that flips of a coin are independent. They are independent given the coin, but when you do not know the coin, earlier flips carry information about it.}
\iqfollow{Is the third flip independent of the first two?}{Not unconditionally, and that is what makes them useful. Given the coin they are independent.}
\end{solution}

\begin{solution}{iq:iv1:conditional-probability-and-bayes:14}
After two positives, the probability is $9/20 = 45\%$. After a negative third test, it is $1/12 \approx 8.3\%$.

Use odds and likelihood ratios\index{likelihood ratio}. Prior odds are $1 : 99$. A positive result has a likelihood ratio of $0.9/0.1 = 9$, and two give $81$, so the odds are $81 : 99 = 9 : 11$, which is $9/20$. A negative test has a ratio of $0.1/0.9 = 1/9$, and that takes the odds to $1 : 11$, which is $1/12$.

The order of the tests makes no difference. In effect, a positive and a negative cancel, and you are left with one net positive on a rare condition.

\iqwrong{Assuming that two positives at 90\% must mean about 99\%. The base rate of 1\% costs you more than two good tests can recover.}
\iqwhy{Multiplying likelihood ratios needs the tests to be independent given the truth. Repeating the same test on the same sample would give the same mistake twice and add little.}
\end{solution}

\begin{solution}{iq:iv1:conditional-probability-and-bayes:15}
Yes, by Simpson's paradox\index{Simpson's paradox}. For example, A wins 9 of 10 trades in the first year and 30 of 100 in the second, while B wins 80 of 100 and then 2 of 10. A is higher in each year (90\% against 80\%, and 30\% against 20\%), yet over the two years A wins 39 of 110 (35\%) and B wins 82 of 110 (75\%).

The cause is the mix. A traded mostly in the hard year and B mostly in the easy year, so the overall figure mainly compares the years. If each traded the same number in each year, the reversal could not happen, because equal-weighted averages preserve the ordering.

Report the rates by year, with the trade counts. If you must give one number, give an average over a common mix of years, because that compares like with like. Which is right depends on the question: is the year a condition you can control, or an outcome?

\iqwrong{Answering ``no'', since a higher rate in every group seems to force a higher rate overall. That holds only with the same weights on each group.}
\iqwhy{Portfolio attribution and performance across regimes run into this all the time. A strategy can look best overall only because it held back in bad periods, or the other way round, and you have to know which.}
\end{solution}

\begin{solution}{iq:iv1:conditional-probability-and-bayes:16}
The probability is $16/35 \approx 0.457$, so a flagged signal is more likely to be false than real.

Count\index{false discovery rate} the outcomes. Fifty signals work, and 40 of them are flagged. Of the 950 that do not work, $47.5$ are flagged. So $40$ of the $87.5$ flagged ones are real.

Tighten the test first. At the 1\% level, keeping the same power for simplicity, the false flags fall to $9.5$ and the true share climbs to about 81\%. In practice the power falls too, so you may need more data, but the base rate is rarely something you can change, and the threshold is. If no signals work at all, every flag is false, and the test gives no warning of that by itself.

\iqwrong{Saying 95\%, since the test is at the 5\% level. That number describes what happens to useless signals, not how likely a flagged signal is to be useful. The answer depends on how many ideas you tried and how good they were.}
\iqwhy{This is the sum behind every research process: with a low base rate of good ideas, most significant results are false unless the test is made much stricter than the usual 5\% level.}
\end{solution}

\begin{solution}{iq:iv1:conditional-probability-and-bayes:17}
The fair price is about $0.567$, or 57 cents.

Each head moves the odds of the $0.6$ coin by a factor of $1.5$ and each tail by $1/1.5$, so only the net count of 4 heads matters\index{likelihood ratio}: the odds are $1.5^4 = 5.0625 : 1$, or a probability of $0.835$ for the $0.6$ coin. The chance of a head next is
\[
0.835 \cdot 0.6 + 0.165 \cdot 0.4 \approx 0.567.
\]
The price must lie between $0.4$ and $0.6$, whatever you have seen, because the only two possibilities are those. Even at very strong evidence, the price approaches $0.6$ from below.

\iqwrong{Quoting $0.7$, the observed frequency. Your beliefs allow only the two coins, and the frequency has no role beyond its likelihood. The data would fit a third coin better, but you are pricing within the beliefs you stated.}
\iqwhy{Prices update through this kind of sum. A market maker quoting the next outcome holds a belief over types, not a frequency, and that is why fair value is a weighted average rather than the latest fill.}
\end{solution}

\begin{solution}{iq:iv1:conditional-probability-and-bayes:18}
The probability is $3/4$.

Update on the box's contents. With $k$ red balls, two reds without replacement has probability $\frac{k(k-1)}{12}$, so the weights on $k = 2, 3, 4$ are $2, 6, 12$, out of 20. The third ball is red with probability $(k - 2)/2$, namely $0$, $1/2$ and $1$. So the answer is $(0 \cdot 2 + \tfrac12 \cdot 6 + 1 \cdot 12)/20 = 3/4$.

This matches Laplace's rule\index{Laplace's rule of succession}, $(2 + 1)/(2 + 2) = 3/4$, and that is no accident. If the number of reds is uniform on $0, \dots, N$, the sequence of draws is exchangeable, and it has the same law as drawing with replacement from a coin with uniform $p$. The size of the box does not matter.

\iqwrong{Guessing a value of $k$ and answering one half or one. It is the uncertainty in $k$ that does the work.}
\iqfollow{And after one red and one blue?}{Then $(1 + 1)/(2 + 2) = 1/2$, by the same rule.}
\end{solution}

\begin{solution}{iq:iv1:conditional-probability-and-bayes:19}
The probability of guilt is about $1/6$.

Start from the population\index{prosecutor's fallacy}. Before the test, this person is the guilty one with probability $1/(5 \times 10^6)$. Besides the guilty person, about 5 of the other 5 million people match by chance. Of about 6 people in the city who match, one is guilty, so a person who was picked with no other evidence and who matches has probability about $1/6$.

By odds: prior odds are about $1 : 5\,000\,000$, and the likelihood ratio is $10^6$, so the posterior odds are $1 : 5$.

The lawyer computes $P(\text{match} \mid \text{innocent})$, and states it as if it were $P(\text{innocent} \mid \text{match})$. The two differ because almost everyone is innocent.

\iqwrong{Treating one in a million as the chance of innocence. Other evidence changes everything, though: with a prior of one half from other evidence, the match makes guilt nearly certain.}
\iqfollow{And if the whole city is tested, and one of the matches is picked at random?}{The count of other matches is then random, about Poisson with mean 5, and the chance is the mean of $1/(1 + K)$, which is $(1 - e^{-5})/5 \approx 0.199$. It is close to $1/6$, but not equal, because the protocol\index{protocol of information} differs.}
\end{solution}

\begin{solution}{iq:iv1:conditional-probability-and-bayes:20}
The probability of skill is about $0.374$, and the chance of beating the market next year is about $0.575$.

Compare the two groups. Skilled managers with five wins make up $0.1 \cdot 0.7^5 \approx 0.0168$ of all managers, and unskilled ones $0.9 \cdot 0.5^5 \approx 0.0281$. So the odds\index{Bayes by odds} are $0.0168 : 0.0281$, and the probability is $0.374$. The forecast for next year is the mix $0.374 \cdot 0.7 + 0.626 \cdot 0.5 \approx 0.575$.

Five straight wins is rarer for a coin flipper, but there are nine times as many coin flippers, and so most managers with this record are still unskilled.

\iqwrong{Treating the record as proof of skill, since $0.5^5$ is only about 3\%. Unskilled managers are numerous, and that 3\% is applied to all of them.}
\iqwhy{Selecting managers on past returns is exactly this calculation, and it says how much a track record can move you: not nothing, but less than the raw numbers suggest.}
\end{solution}

\begin{solution}{iq:iv1:conditional-probability-and-bayes:21}
Bid 0, because every positive bid has a negative expected profit, which is $-b^2/400$. At $2.5\,V$, you should bid 100, for an expected profit of 25.

Condition on the owner's acceptance\index{adverse selection}\index{winner's curse}. A bid of $b$ is accepted exactly when $V \le b$, which has chance $b/100$, and given that, $V$ is uniform on $[0, b]$ with mean $b/2$. So you get something worth $1.5 \cdot b/2 = 0.75b$ on average, and you pay $b$. Each accepted bid loses $0.25b$ on average, and the expected profit is $-(b/100)(0.25b) = -b^2/400$.

At a multiple of 2 the expected value is exactly the bid, and all bids break even. Above that, every bid is profitable, and a higher bid is better. At $m = 2.5$, the profit is $(b/100)(1.25b - b) = b^2/400$, which at the top bid is 25. A bid of 40 at $m = 1.5$ loses 4 on average.

\iqwrong{Bidding on the average value: $V$ averages 50, so it is worth 75 to you, so any bid up to 75 looks profitable. This ignores what you learn from the seller saying yes: the deal only happens when $V$ is low.}
\iqfollow{How large must the multiple be before a positive bid makes sense?}{Above 2. For a multiple $k$, a bid of $b$ earns $\frac{b}{100} \cdot b\bigl(\frac{k}{2} - 1\bigr)$ on average, which is positive only when $k > 2$.}
\iqwhy{Every quote you offer is hit selectively. The same conditioning gives the fill price's true value, and it is why a price that looks good before the trade can look bad straight after it.}
\end{solution}

\begin{solution}{iq:iv1:conditional-probability-and-bayes:22}
The first ten are all red with probability $1/11$. Given a red tenth draw, the first was red with probability $2/3$.

The chance of all red multiplies as $\tfrac12 \cdot \tfrac23 \cdot \tfrac34 \cdots \tfrac{10}{11}$, which telescopes to $1/11$.

The urn looks like it remembers the past, but it is exchangeable\index{exchangeability}\index{P\'olya urn}: any sequence with the same number of reds has the same probability, because the numerators and denominators just appear in a different order. So only the number of reds matters, and in fact the count after $n$ draws is uniform on $0, \dots, n$. So $P(\text{first red} \mid \text{tenth red}) = P(\text{first red} \mid \text{second red}) = \tfrac12 \cdot \tfrac23 \big/ \tfrac12 = \tfrac23$. The later draw says as much about the first as the first says about the later.

\iqwrong{Assuming the future cannot inform the past, and giving $1/2$. Conditioning looks both ways, because both draws depend on the same hidden share of red.}
\iqfollow{What is this urn equal to?}{A coin with a heads probability that is uniform between 0 and 1, and with independent flips given that probability. That is why it matches Laplace's rule.}
\end{solution}

\begin{solution}{iq:iv1:conditional-probability-and-bayes:23}
The flips are not independent. Two heads have probability $1/3$, and the correlation is $1/3$.

Condition on $p$\index{law of total probability}: $P(\text{both heads}) = E[p^2] = \int_0^1 p^2 \, dp = 1/3$. Each flip alone has probability $E[p] = 1/2$, so independence would need $1/4$. They are independent given $p$, but not overall: a head is news about $p$.

The covariance is $E[p^2] - E[p]^2 = \mathrm{Var}(p) = 1/12$, and each flip has variance $1/4$. So the correlation is $(1/12)/(1/4) = 1/3$. A coin with known $p$ would have no correlation at all, since its variance of $p$ is zero. In this sense the correlation measures how much is unknown.

\iqfollow{And with more flips?}{The same correlation holds between any pair, since the flips are exchangeable, and the pairwise correlation does not die out with distance.}
\iqwhy{Returns that look correlated across assets, or across time, may all be independent draws given a hidden state such as a volatility regime. A shared unknown creates dependence with no direct link.}
\end{solution}

\begin{solution}{iq:iv1:conditional-probability-and-bayes:24}
Given $X = 2$, $Y$ has mean 1 and variance $3/4$. And $P(Y > 0 \mid X > 0) = 2/3$.

Write $Y = \rho X + \sqrt{1 - \rho^2}\, Z$ with $Z$ independent of $X$. That is the projection\index{conditional normal} of $Y$ onto $X$. It gives a conditional mean of $\rho x$ and a variance of $1 - \rho^2$, which does not depend on $x$.

For the sign, use the angle in the plane of $(X, Z)$\index{orthant probability}. The events $X > 0$ and $Y > 0$ are half-planes, and their normals have an angle between them whose cosine is $\rho$. The direction of $(X, Z)$ is uniform, so
\[
P(X > 0, Y > 0) = \frac14 + \frac{\arcsin \rho}{2\pi},
\]
and $P(Y > 0 \mid X > 0) = \tfrac12 + \arcsin(\rho)/\pi$. With $\rho = 0.5$ the arcsine is $\pi/6$, which gives $2/3$. As a check, $\rho = 0$ gives $1/2$, and $\rho = 1$ gives 1.

\iqwrong{Answering $0.75$ by reading the $\rho$ as a linear increase in the sign's chance. The relation between sign agreement and correlation is not linear: at $\rho = 0.5$ it is $2/3$.}
\iqwhy{That same formula links sign accuracy and correlation for a signal, so an agreement of $2/3$ means a correlation of one half. Correspondingly, 55\% sign accuracy means a correlation of only about $0.16$.}
\end{solution}

\enlargethispage{-3\baselineskip}%
\begin{solution}{iq:iv1:conditional-probability-and-bayes:25}
The updated mean is $103.84$ and the standard deviation is $2.4$.

For normal beliefs and normal noise, precisions add up\index{precision weighting}, and the mean is the precision-weighted average. The prior has precision $1/16$, the signal $1/9$. So the new precision is $25/144$, which is a variance of $5.76$, or a standard deviation of $2.4$. The signal's weight is $(1/9)/(25/144) = 16/25 = 0.64$, so the mean moves $0.64 \cdot 6 = 3.84$ from 100.

The extreme cases are as you would expect: a perfect signal gets all the weight, and a useless one gets none. And the variance does not depend on what the signal said.

\iqwrong{A simple average, 103, which ignores the precisions, or a final standard deviation near 3 or 4. Combining the two gives a value tighter than either.}
\iqfollow{And after a second, identical signal?}{Precision becomes $1/16 + 2/9$, and the mean moves further. In general, precision simply adds per signal, so long as the noises are independent.}
\iqwhy{This formula underlies the Kalman filter and any shrinkage of a noisy forecast towards a prior.}
\end{solution}

\begin{solution}{iq:iv1:conditional-probability-and-bayes:26}
The lowest safe selling price is 102, and after that sale it should rise to about $103.85$.

The price must equal the expected value of the stock given that someone buys from you\index{adverse selection}. If the value is 110, a buy has probability $0.2 + 0.8 \cdot 0.5 = 0.6$, because all informed traders buy. If it is 90, a buy has probability $0.4$. So $P(110 \mid \text{buy}) = 0.6$, and the expected value given a buy is $0.6 \cdot 110 + 0.4 \cdot 90 = 102$. By the same reasoning the bid is 98, and the spread of 4 pays for the informed traders.

Each trade is news. After a buy, your belief in 110 is $0.6$. A further buy has odds\index{Bayes by odds} $0.6 \cdot 0.6 : 0.4 \cdot 0.4$, or $9/13$, and the ask rises to $90 + 20 \cdot 9/13 \approx 103.85$.

\iqwrong{Quoting a fair price of 100 plus a markup chosen for profit. A quote at 100 loses on average, because buyers are more likely when the stock is high. The spread comes from conditioning on the trade, not from a fee.}
\iqfollow{What are the limits?}{With no informed traders, the ask is 100. With only informed traders, it is 110, and at that price a buyer gains nothing, so no trade is worth making.}
\end{solution}

\begin{solution}{iq:iv1:conditional-probability-and-bayes:27}
The probabilities are $49/58 \approx 0.845$ with independent errors, $0.7$ if they always agree, and $11/14 \approx 0.786$ with the correlated numbers.

Up and down are equally likely, so all you need is the chance of what you saw given each truth. Both say up in a world that goes up exactly when both are right, and in a world that goes down exactly when both are wrong. So $P(\text{up} \mid \text{both say up}) = a/(a + b)$, where $a$ and $b$ are the chances of both right and both wrong\index{correlated evidence}.

Independent errors: $a = 0.49$ and $b = 0.09$. Always agreeing: $a = 0.7$ and $b = 0.3$, and the second analyst adds nothing. Correlated: $a = 0.55$ and $b = 0.15$, giving $11/14$. These numbers are consistent: the other 30\% of days split evenly between the two kinds of disagreement, so each analyst is right on $0.55 + 0.15 = 0.70$ of days, the $0.15$ being the days when that analyst alone is right.

\iqwrong{Multiplying the likelihood ratios for the two analysts, which gives $(7/3)^2$ every time. That is valid only for independent errors. Sources that repeat each other do not add evidence.}
\iqwhy{Combining signals by simple addition of their strength overstates confidence if the signals share their data. That is one reason a portfolio of ten similar signals has far less value than ten independent ones.}
\end{solution}

\begin{solution}{iq:iv1:conditional-probability-and-bayes:28}
Given at least one ace, the probability of at least two is about $0.370$. Given the ace of spades, it is about $0.561$.

Count hands. There are $\binom{52}{13}$ hands, of which $\binom{48}{13}$ have no ace and $4\binom{48}{12}$ have exactly one. So
\[
P(\ge 2 \mid \ge 1) = \frac{\binom{52}{13} - \binom{48}{13} - 4\binom{48}{12}}{\binom{52}{13} - \binom{48}{13}} \approx 0.370.
\]
Given the ace of spades, the other 12 cards come from the remaining 51, and having no other ace has chance $\binom{48}{12}/\binom{51}{12}$. So the chance of at least one more is about $0.561$.

The events differ\index{conditioning on a named item} because ``holds the ace of spades'' is more specific. A hand with two aces fits ``holds the ace of spades'' half of the time, while a hand with one ace fits only a quarter of the time. So naming the card favours hands with more aces, exactly as in the case of the boy born on a Tuesday.

\iqwrong{Assuming that the suit cannot matter, because aces are symmetric. Each suit has the same chance, but a named suit is less likely to be held, and more so for hands with fewer aces.}
\iqfollow{What if a player chooses to name their ace's suit when they hold one?}{Then the suit is no extra information about the number of aces, and $0.370$ stands. The protocol\index{protocol of information} is what decides it.}
\end{solution}

\begin{solution}{iq:iv1:conditional-probability-and-bayes:29}
Keep your door at the first opening, and switch at the second. You win with probability $3/4$.

There are four plans. Keep and keep wins $1/4$, because the host's actions tell you nothing about your own door. Keep and switch wins whenever your first pick was wrong, since after two openings the only closed door that is not yours must hold the prize: $3/4$. Switching first and then keeping wins $3/8$. Switching both times wins $5/8$.

The rule is to hold on as long as you can, then switch at the last moment. At each step the host is forced to give you information about the other doors, while your door's probability is stuck at its initial value of $1/4$. So you want to hold the door with the lowest probability for as long as possible, and switch only when the switch collects all the rest\index{protocol of information}.

To check the $5/8$ for switching twice: your first switch has chance $3/8$ of holding the prize, the host opens one more door, and you switch to the remaining door. You win unless you held the prize at the point of the second switch, which leaves $5/8$. Exact enumeration over all paths and all the host's choices confirms each of the four numbers, and a simulation matches the best one.

\iqwrong{Switching as often as possible, because switching once was good in the three-door game. Each early switch leaves the door that the host cannot touch, which is exactly the door whose probability you want to transfer later.}
\end{solution}

\begin{solution}{iq:iv1:conditional-probability-and-bayes:30}
A family has 2 children on average, half of all children are boys, and a randomly met child comes from a family of 3 children on average.

The family size $L$ is geometric on $1, 2, \dots$ with $P(L = k) = 2^{-k}$, so $E[L] = 2$. Each family has exactly one boy, so boys are one child in two: the rule does not change the sex ratio, since every birth is an independent coin.

Picking a child is different from picking a family. A family with $k$ children is $k$ times as likely to be hit, so this is a size-biased sample\index{inspection paradox}\index{size-biased sampling}:
\[
\frac{E[L^2]}{E[L]} = \frac{\mathrm{Var}(L) + E[L]^2}{E[L]} = \frac{2 + 4}{2} = 3.
\]
This is the same effect as the longer bus gap that a random arrival falls into. If all families had exactly 2 children, the two averages would match.

\iqwrong{Saying that the stopping rule favours boys, or that a random child's family averages 2. A stopping rule can never bias a fair coin\index{optional stopping}, but it does change how you see the families when you sample them by child.}
\iqwhy{Whenever you sample per item rather than per group, the larger groups are overweighted. This is why the average investor's fund is larger than the average fund, and why the average trade occurs at a busier time than the average minute.}
\end{solution}

\begin{solution}{iq:iv1:conditional-probability-and-bayes:31}
The correlation is $-1/2$, even though the scores are independent in the population.

Those hired are uniform on the triangle above the line $x + y = 1$. On this triangle, $E[X] = 2/3$, $E[X^2] = 1/2$, so $\mathrm{Var}(X) = 1/18$. Integrating gives $E[XY] = 5/12$, so the covariance is $5/12 - 4/9 = -1/36$, and the correlation is $(-1/36)/(1/18) = -1/2$.

The reason is that a person with a low $X$ can be hired only with a high $Y$. This selection effect\index{Berkson's bias}\index{selection bias} appears with any hurdle on the sum. It grows as the hurdle rises to 1, and then stays at exactly $-1/2$, because above 1 those hired form a smaller copy of the same triangle.

So studying only hires can make it seem that the two skills trade off, or that one of them predicts poor performance. It is a property of the filter, not of people.

\iqwrong{Assuming that a selection keeps independence. Filtering on a function of both variables creates dependence where there was none.}
\iqwhy{The same trap appears when a study uses only funds that survived, or trades that were filled. Any filter on outcome can produce a correlation between two inputs that have nothing to do with each other.}
\end{solution}

\begin{solution}{iq:iv1:conditional-probability-and-bayes:32}
After $-4\%$, the probability of crisis is about $0.956$. After a return of 0, it falls to about $0.017$.

The update needs the density ratio\index{likelihood ratio} at the observed return:
\[
\frac{f_{\text{crisis}}(r)}{f_{\text{calm}}(r)} = \frac{1/3 \cdot e^{-r^2/18}}{e^{-r^2/2}} = \frac13 e^{\,4r^2/9}.
\]
At $r = 4$, with returns in percent, this is $\tfrac13 e^{7.11} \approx 409$. The prior odds are $5 : 95$, so the posterior odds are about $21.5 : 1$, which gives $0.956$. So one single move can do almost all the work, since a move of 4 standard deviations is exceptionally rare in the calm regime.

At $r = 0$, the ratio is $1/3$, since a return near zero is three times as likely in the calm regime, where the density is more peaked. The odds become $5 : 285$, or about $0.017$. The two densities are equal at $|r| \approx 1.57\%$, and returns smaller than that actually lower your belief in a crisis.

\iqwrong{Treating $-4\%$ as ``about 1.3 standard deviations under the crisis'', and so not strong evidence. The evidence is the comparison between two models, and a move that is unremarkable under one and extreme under the other is very strong.}
\iqwhy{This is one step of the filter behind regime models and of GARCH's behaviour. Quiet days should lower your risk estimate, and one large move should raise it fast.}
\end{solution}

\begin{solution}{iq:iv1:conditional-probability-and-bayes:33}
For exponentials, $X$ given $X + Y = s$ is uniform on $[0, s]$. For standard normals, it is normal with mean $s/2$ and variance $1/2$.

Exponentials: the joint density at $(x, s - x)$ is $e^{-x} e^{-(s - x)} = e^{-s}$, and this does not depend on $x$. So every split of $s$ is equally likely\index{conditioning on a sum}: the mean is $s/2$ and the variance is $s^2/12$.

Normals: the density is proportional to $e^{-x^2/2 - (s-x)^2/2}$, which by completing the square is a normal in $x$ with mean $s/2$ and variance $1/2$. The variance does not depend on $s$ at all.

For any independent and identically distributed pair, the mean of $X$ given the sum is $s/2$ by symmetry, so that is no test of which law you have. The spread is what tells them apart.

\iqwrong{Expecting $X$ to be close to $s/2$ in every case. For exponentials, a large sum is equally likely to come from any split, including one big value and one tiny one. For normals, a big sum is shared out.}
\iqwhy{This distinction matters for tail risk. With fat-tailed returns, a large move over two days most likely came from one extreme day. With normal returns, it came from two moderately large days. Your conclusions about the path depend on which model you hold.}
\end{solution}

\begin{solution}{iq:iv1:conditional-probability-and-bayes:34}
The error is in treating ``the other envelope holds $2x$ or $x/2$ with equal chance'' as true for every $x$. There is no uniform prior over all amounts. Given a proper prior, this fails for some amounts, and those exactly cancel the gain.

Here is the example. The six equally likely cases are: you hold 1 or 2 from the first pair, 2 or 4 from the second, and 4 or 8 from the third. So you hold 1 and 8 with chance $1/6$ each, and 2 and 4 with chance $1/3$ each.
\begin{itemize}
\item Holding 1, the other holds 2: a gain of 1.
\item Holding 2, the other holds 1 or 4: an average of $2.5$, a gain of $0.5$.
\item Holding 4, the other holds 2 or 8: an average of 5, a gain of 1.
\item Holding 8, the other holds 4: a loss of 4.
\end{itemize}
The average gain\index{two-envelope paradox} is $\tfrac16(1) + \tfrac13(0.5) + \tfrac13(1) + \tfrac16(-4) = 0$. Both envelopes average $3.5$. So switching wins for almost every amount you might see, and the big loss at the top exactly pays for those gains.

\iqwrong{``Since it holds for each $x$, you should switch.'' It does not hold for each $x$, as the top value shows. If you tried to remove that top value, you would need an unbounded prior, and then the expected amount is infinite, which is where the paradox lives.}
\end{solution}

\begin{solution}{iq:iv1:conditional-probability-and-bayes:35}
The probability of the 30-minute timetable is $2/3$, and you will wait about $7.5$ more minutes on average.

A random arrival waits uniformly up to the gap. Waiting 10 minutes with no bus has chance $5/15 = 1/3$ under the 15-minute timetable and $20/30 = 2/3$ under the 30-minute one. So the odds\index{Bayes by odds} are $2 : 1$ for the 30-minute timetable.

Given the 15-minute timetable and no bus yet, the remaining wait is uniform on $[0, 5]$, for a mean of $2.5$. Given the 30-minute one, it is uniform on $[0, 20]$, for a mean of 10. So the expected remaining wait is $\tfrac23 \cdot 10 + \tfrac13 \cdot 2.5 = 7.5$ minutes. Before you waited, the expected wait was $\tfrac12(7.5) + \tfrac12(15) = 11.25$, so after 10 minutes the remaining wait is $7.5$, not $1.25$: waiting has told you that the buses are less frequent than you hoped.

\iqwrong{Using the average of the two timetables without updating, or assuming that you must be close to the bus. Having waited is evidence about the timetable, and that keeps the expected wait high.}
\iqwhy{Any estimate of ``time to next'' behaves like this. If an expected fill or a recovery does not come, then the right response is often to lengthen the estimate, not to feel that it must be close.}
\end{solution}

\begin{solution}{iq:iv1:conditional-probability-and-bayes:36}
The chance that everyone settles on the wrong type is $1/5$, and on average $3.6$ people guess before the first person whose signal no longer matters.

Guess by guess, reduce the state\index{information cascade}: let $d$ be the number of earlier guesses for one type minus those for the other. While $|d| < 2$, each person's guess reveals their signal. With $d = 0$ they follow their own signal. With $d = 1$ a signal that agrees adds a guess, and a signal that disagrees gives a tie, which by the rule they follow. But at $d = 2$ even a disagreeing signal leaves a net one signal in favour of the crowd. So they ignore their own signal. Their guess reveals nothing, the next person faces the same state, and the crowd is stuck.

So $d$ is a random walk, moving towards the truth with probability $p = 2/3$ and away with $q = 1/3$, absorbed at $\pm 2$. Moving in two-step pairs\index{gambler's ruin}, the probability of being absorbed at the wrong end is $q^2/(p^2 + q^2) = 1/5$, and the average time to absorption is $2/(p^2 + q^2) = 18/5$.

Simulating people who do the full Bayesian update on public guesses gives the same numbers.

\iqwrong{Thinking that a crowd must eventually be right because signals are good and there are many people. Once a cascade starts, no new information reaches the public, even though almost all the private signals remain unused.}
\iqwhy{Prices can stop revealing information when traders copy each other's orders, and a price held by a cascade can break on small news.}
\end{solution}

\begin{solution}{iq:iv1:conditional-probability-and-bayes:37}
The best expected profit is $\$0.08$. Knowing the coin's type would give $\$0.60$.

Do it by rounds, always betting on the likelier side given what you have seen\index{sequential update}. With even-money \$1 bets, a bet with probability $u$ of winning earns $2u - 1$.

Round 1: either side has $u = 1/2$, so it earns 0. Round 2: after one flip, the belief in the favoured coin is $0.6$, and $u = 0.6 \cdot 0.6 + 0.4 \cdot 0.4 = 0.52$, so it earns $0.04$. Round 3: after HH or TT, which has total probability $0.52$, the odds are $9 : 4$, and $u = 7/13$, so the bet earns $1/13$. After HT or TH, the belief returns to even, and the bet earns 0. The expected earning for round 3 is $0.52/13 = 0.04$. In total, this is $0.08$.

Knowing the coin, you earn $0.2$ a round, so three rounds give $0.6$. Most of that edge is out of reach, because you must first learn which coin you hold. Exact enumeration and a simulation of the rule both give $0.08$.

\iqwrong{Treating the coin as fair on average and answering zero. The prior average is fair, but you learn as you go, and betting on what you learn has positive value, though it is small early on.}
\iqwhy{The early rounds are worth almost nothing except as information. That is why a trader might trade small for the sake of learning, and why you must count that value when judging whether a first trade makes sense.}
\end{solution}

\begin{solution}{iq:iv1:conditional-probability-and-bayes:38}
With 60 of 100, the probability of fair is about $0.523$. With $5\,100$ of $10\,000$, it is about $0.915$, even though the p-value is about the same. This is Lindley's paradox\index{Lindley's paradox}.

The chance of 60 heads under a fair coin is $\binom{100}{60}2^{-100} \approx 0.0108$. Under the uniform prior, every count from 0 to 100 is equally likely, so that chance is $1/101 \approx 0.0099$. So the data are a little more likely under fairness, which gives $0.523$.

With 10,000 flips, the fair chance is about $0.00108$ (the normal density at $z = 2$ divided by 50), and the flat chance is $0.0001$. Now the data favour fairness 11 to 1.

The two-sided p-values are $0.057$ and $0.047$. Those measure how surprising the count is under the fair coin. The Bayes factor compares two explanations, and a uniform alternative spreads its chance over all values, most of which are nowhere near $0.51$. So the data are ``unlikely under fair, but less likely still under the vague alternative''.

The conclusion depends strongly on the alternative. Against the single alternative $p = 0.6$, the likelihood ratio is about $7.5$ in favour of $p = 0.6$, and the probability of fair falls to about $0.12$.

\iqwrong{Saying that a p-value of 5\% means a 95\% chance of a real effect, or a 5\% chance of the null. It means neither of those, and the chance of the null can even rise with more data at the same p-value.}
\iqwhy{Large samples make any small effect significant. For a strategy with a vast sample and a borderline t-statistic, ask for the size of effect it implies, and not only the p-value.}
\end{solution}

\begin{solution}{iq:iv1:conditional-probability-and-bayes:39}
The best shows about $1.77$, and its expected true Sharpe ratio is about $0.89$. With no skill at all, the best still shows about $1.25$, and its expected true Sharpe ratio is 0.

Each measured value is true Sharpe plus noise, with total variance $0.25 + 0.25 = 0.5$. The expected maximum of 100 standard normals\index{expected maximum of normals} is about $2.51$, so the best shows $2.51\sqrt{0.5} \approx 1.77$.

Now shrink. For normals, the mean of the true value given the measurement is $\tfrac{0.25}{0.25 + 0.25} = \tfrac12$ of the measurement\index{shrinkage}, whichever manager you pick. Selecting the best changes nothing in that rule, since the rule holds for each value, so the best one's expected true Sharpe ratio is $\tfrac12 \cdot 1.77 \approx 0.89$. Half of its apparent quality is luck.

With no skill, the noise is all there is: the best shows $2.51 \cdot 0.5 \approx 1.25$, and the true value is exactly 0.

\iqwrong{Taking the best one's figure at face value, or subtracting an adjustment for the number of tries. The right correction depends on how much real variation exists, and that needs an estimate of the prior, which you can get from the spread across all 100 managers.}
\iqfollow{How would you estimate the prior's spread?}{The measured variance across managers is the true variance plus the noise variance, so subtract the noise. A measured standard deviation near $0.7$ suggests a true spread near $0.5$.}
\iqwhy{Performance fades after selection for this reason, without any change in skill. Expect half of a backtest's edge to vanish at this level of noise, however many strategies you tried; more if the noise is larger against the true spread.}
\end{solution}

\begin{solution}{iq:iv1:conditional-probability-and-bayes:40}
Yes. Choose a random threshold from any distribution that gives positive probability to every interval, and say ``larger'' if the number you saw is below your threshold. For 10 and 20 with that threshold, your chance is about $0.691$.

Call the numbers $a < b$, and let $T$ be your threshold\index{random threshold strategy}. There are three cases. If $T < a$, you always say ``smaller'', and that is right exactly when you saw $b$: a chance of $1/2$. If $T > b$, you always say ``larger'', and that is right exactly when you saw $a$: also $1/2$. If $a < T < b$, you say the right thing whichever card you turned. So your chance of being right is
\[
\tfrac12 + \tfrac12\, P(a < T < b),
\]
which is strictly above $\tfrac12$ as long as $T$ has positive probability on every interval. For 10 and 20, with $T$ normal with mean 15 and standard deviation 10, $P(10 < T < 20) \approx 0.383$, so you get $0.691$.

The guarantee cannot be made uniform. An opponent who knows your rule can put $a$ and $b$ very close together, or very far from where your threshold lies, and the excess can be made as small as they like. It is positive for each pair, but it has no positive lower bound.

\iqwrong{``You know nothing about how the numbers were chosen, so you can do no better than chance.'' This is true of a fixed threshold, because an opponent can place both numbers on one side of it. Randomisation is what removes that weakness.}
\iqwhy{Having no information about the size of a number does not mean that the observation is useless. A value relative to any reference, however crude, tells you something about the ordering, and that gives a small edge.}
\end{solution}
